\documentclass[tikz, border=10pt]{standalone}

\begin{document}
\begin{tikzpicture}
  % CONFIGURACIÓN
  % Define el área del billete/fondo
  \clip (0,0) rectangle (12,7);
  \fill[white] (0,0) rectangle (12,7); % Fondo base

  % --- CAPA 1: ONDAS DE FONDO (Cyan) ---
  % Iteramos desde abajo (0) hasta arriba (7) en pasos muy pequeños (0.05)
  \foreach \y in {0, 0.05, ..., 7.5} {

    % EXPLICACIÓN DE LA FUNCIÓN:
    % (\x, {\y + ...}) -> La línea base está en 'y'
    % sin(150*\x)      -> Frecuencia horizontal (que tan rápido vibra)
    % + 50*\y          -> DESPLAZAMIENTO DE FASE (La magia: la curva
    % cambia según la altura)

    \draw[cyan!60!black, opacity=0.4, ultra thin]
    plot[domain=0:12, samples=150, variable=\x]
    (\x, { \y + 0.15 * sin(180*\x + 60*\y) });
  }

  % --- CAPA 2: INTERFERENCIA (Magenta - Opcional para efecto Moiré) ---
  % Hacemos lo mismo pero con el ángulo invertido o diferente frecuencia
  % Esto crea esos "diamantes" visuales típicos de los billetes.
  \foreach \y in {0, 0.05, ..., 7.5} {
    \draw[magenta!60!gray, opacity=0.3, ultra thin]
    plot[domain=0:12, samples=150, variable=\x]
    (\x, { \y + 0.15 * sin(140*\x - 40*\y) });
  }

  % Marco simple para delimitar
  \draw[line width=1pt, color=blue!30!black] (0,0) rectangle (12,7);

\end{tikzpicture}
\end{document}

